\lecture{22}{May 26}{}
We now finish the proof of the lemma in the proof of Vizing.
\begin{lemma}[Schrijver]
    If \(\exists v \in V(G)\), and \(k \in \mathbb{N}\) s.t. 
    \begin{itemize}
        \item \(d(v) \le k\)
        \item For any \(u \in N(v)\), \(d(u) \le k\) 
        \item \(d(u) = k\) for at most one neighbor \(u \in N(v)\),     
    \end{itemize}  
    then 
    \(G - v\) is \(k\)-edge colorable iff \(G\) is \(k\)-edge colorable.     
\end{lemma}

\begin{proof}
    We can do induction on \(k\). 

    \textbf{Base case:} \(k = 1\), which is trivial.

    \textbf{Induction step:} We may assume 
    \begin{itemize}
        \item \(d(v) = k\)
        \item \(d(u) = k\) for one neighbor \(u\) of \(v\)  
        \item \(d(u) = k - 1\) for all other neighbor of \(v\) 
    \end{itemize} 
    since, if needed, we can increase the degrees by adding edges to new vertices, and we can color these edges with unused colors (no restriction from new vertices). Suppose we have a proper \(k\)-edge-coloring of \(G - v\). For each color \(i \in [k]\), let \(X_i \subseteq N(v)\) be the set of neighbors where the color \(i\) is unused. 

    Choose a \(k\)-edge-coloring \(\varphi \) of \(G - v\) that minimizes \(\sum_{i=1}^k \vert X_i \vert^2  \).  

    \textbf{Case 1:} \(\exists i\) s.t. \(\vert X_i \vert = 1\), \(X_i = \left\{ u \right\}  \). 

    Let \(G^{\prime} = G - \left\{ u, v \right\} - \varphi ^{-1}(i) \), i.e. we remove all \(i\)-colored edges, as well as \(\left\{ u, v \right\} \). Left with a \((k - 1)\)-edge-coloring of \(G^{\prime} - v\), and \(d_{G^{\prime}} (v) = k - 1\), \(d_{G^{\prime} } (u) = d_G(u) - 1\) for all \(u \in N(v)\). By induction, \((G^{\prime} - v) \) is \((k - 1)\)-edge-colorable, so \(G^{\prime} \) can be \((k - 1)\)-edge-colorable. Now add back the edges in \(G - G^{\prime}\), all of the same color, then we get a proper \(k\)-edge-coloring of \(G\). 

    \textbf{Case 2:} \(\nexists i\) s.t. \(\vert X_i \vert = 1 \). 

    Observe that 
    \begin{align*}
        \sum_{i=1}^k \vert X_i \vert &= \sum_{u \in N(v)} \left( k - (d(u) - 1) \right) = \left( k - (k - 1) \right) + (k - 1) (k - (k - 1 - 1)) = 2k - 1.     
    \end{align*}     
    Note that the first equality holds since we can count on every neighbor \(u\) of \(v\) and for each \(u\), it has \(k - (d(u) - 1)\) colors unused since every edge incident to \(u\) except \(\left\{ u, v \right\} \) has a color, distinct from those of all the other edges incident to \(u\). 

    Since \(\vert X_i \vert \neq 1 \) for any \(i\) and \(\sum_{i=1}^k \vert X_i \vert = 2k - 1  \), there must be some \(i\) s.t. \(\vert X_i \vert = 0 \), otherwise the sum is odd, there must be some \(\vert X_j \vert \ge 3 \), and \(\vert X_j \vert \ge 2 \) for all \(j\), which implies the sum must be larger than \(2k + 1\).

    Let \(u \in N(v)\) s.t. \(u \in X_j\). Consider the component of \(u\) in the \((i \, \& \, j)\)-colored subgraph. (a path, since it is the union of two matchings.) Swap the colors \(i\) and \(j\) in this component. Let \(X_{\ell }^{\prime} \) be the \(X_{\ell }\) for this new coloring. Then 
    \[
        \vert X_{\ell }^{\prime}  \vert = \begin{dcases}
            \vert X_{\ell} \vert , &\text{ if } \ell \neq i, j ;\\
            1 \text{ or } 2 , &\text{ if } \ell = i ;\\
            \vert X_j \vert - 1 \text{ or } \vert X_j \vert - 2   , &\text{ if } \ell = j,
        \end{dcases}
    \]              
    and we still have a proper \(k\)-edge-coloring of \(G - v\). Since \(X_i^{\prime} \) and \(X_j^{\prime} \) are much closer to each other, and other \(X_{\ell ^{\prime} } \) remain unchanged, we have \(\sum_{i} \vert X_i^{\prime}  \vert^2 < \sum_{i} \vert X_i \vert^2    \), contradicting the choice of \(\varphi \).                             
\end{proof}

\begin{definition}
    The list-chromatic index \(\chi _\ell ^{\prime} (G)\) is the smallest \(k\) s.t. whenever each edge \(e\) of \(G\) is assigned a list \(L_e\) of at least \(k\) colors, there is a proper edge-coloring where every edge uses a color from its list.     
\end{definition}

\begin{remark}
    \(\chi _\ell ^{\prime} (G) = \chi _\ell (L(G))\). 
\end{remark}

Recall that for any general graph \(G\), \(\chi _\ell (G)\) can be arbitrarily larger than \(\chi (G)\) since if \(m \ge \binom{2k - 1}{k}\), \(\chi _\ell (K_{m, m}) \ge k\).

\begin{conjecture}[List-coloring conjecture]
    For any graph \(G\), \(\chi _\ell ^{\prime} (G) = \chi ^{\prime} (G)\).  
\end{conjecture}

\begin{observation}
    \begin{align*}
        \chi ^{\prime} (G) &\le \chi _\ell ^{\prime} (G) \quad \text{[give every edges the same list of colors]}  \\
        &= \chi _\ell (L(G)) \\
        &\le \Delta (L(G)) + 1  \quad \text{[greedy]} \\
        &\le 2 \left( \Delta (G) - 1 \right) + 1 \\
        &\le 2 \chi ^{\prime} (G) - 1 \quad [\chi ^{\prime} (G) \ge \Delta (G)] . 
    \end{align*}
\end{observation}

\begin{theorem}[Kahn, 1996]
    As \(\Delta(G) \to \infty \), 
    \[
        \chi _\ell ^{\prime} (G) = (1 + o(1)) \chi ^{\prime} (G).
    \] 
\end{theorem}

\begin{theorem}[Galvin, 1995]
    If \(G\) is bipartite, then
    \[
        \chi _\ell ^{\prime} (G) = \chi ^{\prime} (G).
    \] 
\end{theorem}

\begin{proof}
    A special case is \(G = K_{n, n}\), which is called the Dinitz Conjecture. 

    We will 
    \begin{enumerate}
        \item Prove the general upper bound in list-chromatic \textbf{number} (vertex coloring) 
        \item Apply this to \(L(K_{n, n})\). 
    \end{enumerate}

    \begin{definition}
        In a directed graph \(\vv{D}\), a kernel is a set of vertices \(I \subseteq V(\vv{D})\) such that 
        \begin{itemize}
            \item \(I\) is independent (no edges in \(I\))
            \item for every \(v \notin I\), \(\exists u \in I\) s.t. \((v, u) \in E(\vv{D}) \).    
        \end{itemize}  
    \end{definition}

    \begin{remark}
        Not every \(\vv{D}\) has a kernel. For example, a directed cycle of 5 vertices. 
    \end{remark}

    \begin{definition}
        A directed graph is kernel-perfect if for every subset \(S \subseteq V(\vv{D})\), the induced subgraph \(\vv{D} [S]\) has a kernel.  
    \end{definition}

    \begin{observation}
        Every (undirected) graph \(G\) admits a kernel-perfect orientation. This is because we can order the vertices arbitrarily: \(v_1, \dots , v_n\), then orient edges from right to left: if \(i < j\), then \(\left\{ v_i, v_j \right\} \to (v_j, v_i)\). Now give any \(S \subseteq V(G)\), run greedy coloring algorithm on \(G[S]\), and set of color 1 vertices is a kernel.     
    \end{observation}

    \begin{lemma}[Bopanna-Bondy-Seidel]
        Let \(G\) be an undirected graph, and let \(\vv{D}\) be a kernel-perfect orientation of \(G\). Then 
        \[
            \chi _\ell (G) \le \Delta ^+ (\vv{D}) + 1,
        \]   
        where \(\Delta ^+(\vv{D})\) is the maximum number of out degree. 
    \end{lemma}

    \begin{proof}
        We have 
        \begin{itemize}
            \item a kernel-perfect orientation \(\vv{D}\) of \(G\). 
            \item for every vertex \(v\), a list \(L_v\) of \(\Delta ^+(\vv{D}) + 1\) colors.      
        \end{itemize}
        Need to assign each vertex a color from this list. Equivalently, decide for each color, which vertices receive that color. 

        Now let \(c \in \bigcup_{v} L_v \) be a color, and let 
        \[
            S_c = \left\{ v \in V(\vv{D}): c \in L_v \right\} = \text{ all vertices that have } c \text{ on their list}.  
        \]  
        Also, let \(I_c \subseteq S_c\) be a kernel of \(\vv{D} [S_c]\). Assign \(c\) to every vertex in \(I_c\). Since \(I_c\) is independent, such coloring is proper. We do this to every color \(c\) (if some vertices already colored, then don't color it again).  
        
        Now we show that every vertex is assigned a color after doing above step. If \(c \in L_v\) and \(v \notin I_c\), for some \(c \in L_v\), then since \(I_c\) is a kernel, there is some vertex \(u \in I_c\) s.t. \((v, u) \in E\). Thus, for any vertex \(v\) there are at most \(\Delta ^+ (\vv{D}) \) such vertices \(u\), i.e. there are at most \(\Delta ^+ (\vv{D})\) colors \(c\) s.t. \(c \in L_v\) but \(v \notin I_c\), but \(\vert L_v \vert > \Delta ^+(\vv{D}) \), so there is some color \(c \in L_v\) with \(v \in I_c\).

        Thus, we have a proper list coloring.              
    \end{proof}

    Therefore, to prove Galvin's theorem, it suffices to find a kernel-perfect orientation \(\vv{D}\) of \(L(K_{n, n})\) s.t. \(\Delta ^+(\vv{D}) = n - 1\). 
\end{proof}